\documentclass[10pt]{article}

% Specify the margins and text size.
\setlength{\textwidth}{6.4in}
\setlength{\textheight}{9.5in}
\setlength{\oddsidemargin}{0pt}
\setlength{\evensidemargin}{0pt}
\setlength{\topmargin}{0pt}
\setlength{\hoffset}{.05in}
\setlength{\voffset}{-1in}

\setlength{\parskip}{5pt}
\setlength{\parindent}{0pt}

% Load some fonts and symbol packages
\usepackage{latexsym}
\usepackage{pifont}       % contains 'star' symbol for counterinsurgency handbook title
%\usepackage{yfonts} 
\usepackage{amsmath}
\usepackage{amsfonts}

\usepackage{graphicx}     % actually, this is loaded by pstricks
\usepackage[T1]{fontenc}
\usepackage{ifthen}
\usepackage{pstricks,pst-grad,pst-text,pst-node,multido,pst-plot,calc,pst-3dplot}
%\usepackage[all]{xy}
%\usepackage{animate}

% The hyperref package inserts links into the pdf.
\definecolor{MyLinkColor}{rgb}{.1,.2,1}
\definecolor{MyCiteColor}{rgb}{.1,1,.2}
\definecolor{MyURLColor}{rgb}{.4,.4,.4}
\usepackage[backref=true,pagebackref=false,hyperindex,colorlinks=true,
  linkcolor=MyLinkColor,urlcolor=MyURLColor]{hyperref}


% The tweaklist package is something I found on the web.  It provides a simple interface
% for making changes to spacing used in the itemize and enumerate environments.  Comment
% this out if you don't care to use tweaklist.
\usepackage{tweaklist}
\renewcommand{\itemhook}{\setlength{\parskip}{2pt}\setlength{\parsep}%
{1pt}\setlength{\topsep}{0pt}\setlength{\itemsep}{0pt}}

\newcommand{\U}{\underline{\hspace{5pt}}}

\usepackage{listings}
\newcommand{\Z}{\hphantom{0}}

\begin{document}
\pagestyle{empty}
\lstset{language=R, showspaces=false, showstringspaces=false}
\href{http://www.su.edu}{\includegraphics[height=1.75cm]{sulogo.eps}}
\vspace{-1.79cm}

{{\ }\hfill\small
\begin{tabular}{cl}
& Math 207\\
& Introduction to Statistics\\
\end{tabular}
}

\setlength{\baselineskip}{1.1\baselineskip}
\bigskip

\newcommand{\BL}{\underbar{\hspace{.5in}}}
\begin{center}
\textbf{\large Review Exercises}
\end{center}
1. For entering freshmen at a certain university, scores on the 
Math SAT and Verbal SAT can be summarized as follows.\vspace{-4pt}
\begin{center}
\begin{tabular}{lll}
$\mbox{average M-SAT} = 555$, & $\mbox{SD}=125$, & \\
$\mbox{average V-SAT} = 543$, & $\mbox{SD}=115$, & $r=0.66$\vspace{-4pt}
\end{tabular}
\end{center}
The scatter diagram is football-shaped.  One student is chosen at random and has
an M-SAT score of 600. You would guess his V-SAT is \BL\ points,
and would have about a 68\% chance to be right within \BL\ points.
\vspace{1.25in}


2. A die is rolled 6 times.  Find the chance that the first number
rolled comes up three more times ---

\hspace{20pt} a) If the first roll is an ace.
\vspace{.7in}

\hspace{20pt} b) If the first roll is an six.
\vspace{.7in}

\hspace{20pt} c) If you don't know what happens on the first roll.
\vspace{1in}

3. A die will be rolled 20 times.  The sum
\[\mbox{number of ones rolled} + \mbox{number of sixes rolled}\]
will be around \BL, give or take \BL\ or so.

\vfill
\eject
{\ }

4. A box contains 10,000 marbles:  6,000 are red and 4,000 
are blue;  500 marbles are drawn at random without replacement.

\hspace{20pt} a) Suppose that there are 218 blue marbles in the sample.
Find the expected value for the percentage of blues in the sample, the observed value, the
chance error, and the standard error.
\vspace{1.5in}

\hspace{20pt} b) Suppose there are 191 blue marbles in the sample.  Repeat a).
\vspace{1.5in}

5.
The town of Hayward California has about 5,000 registered voters.
A political scientist takes a simple random sample of 500 of these voters.
In the sample, the breakdown by party affiliation is:
\begin{center}
\begin{tabular}{lr}
Republican & 115\\
Democrat & 331\\
Independent & 54
\end{tabular}
\end{center}

\hspace{20pt} a) Among all registered voters in Hayward, the percentage of independents is 
estimated as \BL.
\vspace{.5in}

\hspace{20pt} b) This estimate is likely to be off by \BL\ or so.
\vspace{.5in}

\hspace{20pt} c) The range from \BL\ to \BL\ is a 95\% confidence interval for the 
percentage of independents.
\vspace{.5in}

\hspace{20pt} d) Repeat a)--c) for Democrats.

\vfill
\eject
\href{http://www.su.edu}{\includegraphics[height=1.75cm]{sulogo.eps}}
\vspace{-1.79cm}

{{\ }\hfill\small
\begin{tabular}{cl}
& Math 207\\
& Introduction to Statistics\\
\end{tabular}
}
\bigskip

\begin{center}
\textbf{\large  Confidence Intervals (III)}
\end{center}

1. (Hypothtical)  The British Imperial Yard is sent to Paris for calibration
against The Meter.  The length is determined 100 times.  This sequence of
measurements averages out to 91,4402 cm, and the SD is 800 microns (a \textit{micron}
is a millionth of a meter).

\hspace{20pt} a) Is a single reading off by around 80 microns, or 800 microns?
\bigskip

\hspace{20pt} b) Is the average of all 100 readings off by around 80 microns, or 800 microns?
\bigskip

\hspace{20pt} c) Find a 95\%-confidence interval for the exact length of the
Imperial Yard.
\vspace{.5in}

2. 
The Bureau is about to weigh a one-kilogram checkweight 100 times, and take
the average of the measurements.  They are willing to assume the
Gauss model, with no bias, and on the basis of past experience they 
estimate the SD of the error box to be 50 micrograms.

\hspace{20pt} a) The average of all 100 measurements is likely to be off
the exact weight by \BL\ or so.
\bigskip

\hspace{20pt} b) The SD of all 100 measurements is likely to be around \BL.
\bigskip
\bigskip

\hspace{20pt} c) Esimate the probability that the average of all 100 measurements
will be within 10 micrograms of the exact weight.
\vspace{.7in}

3. 
Laser altimeters can measure elevation to within a few inches, without bias,
and with no trend or pattern to the measurements.
As part of an experiment, 25 readings were made on the elevation of a mountain peak.
These averaged out to 81,411 inches, and their SD was 30 inches. 

\hspace{20pt} a) The elevation of the mountain peak is estimated 
as \BL; this estimate is likely to be off by\vspace{-4pt}

\hspace{20pt}\hphantom{f) } \BL\ or so.
\bigskip
\bigskip

\hspace{20pt} b) (T/F) 81,411$\pm$ 12 inches is a 95\%-confidence interval for the elevation
of the mountain peak. 
\bigskip
\bigskip

\hspace{20pt} c) (T/F) 81,411$\pm$ 12 inches is a 95\%-confidence interval for 
average of the 25 readings.

\hspace{20pt} d) (T/F) There is about a 95\% chance that the next reading will be in 
the range 81,411 $\pm$ 12 inches.

\hspace{20pt} e) (T/F) About 95\% of the readings were in the range 81,411 $\pm$ 12 inches.

\hspace{20pt} f) If another 25 readings are made, there is about a 95\% chance that their
average will be in the range\vspace{-4pt}

\hspace{20pt}\hphantom{f) } 81,411 $\pm$ 12 inches.
\vfill
\eject


\end{document}

